\documentclass[conference]{IEEEtran}
\IEEEoverridecommandlockouts
\usepackage{cite}
\usepackage{amsmath,amssymb,amsfonts}
\usepackage{algorithmic}
\usepackage{graphicx}
\usepackage{textcomp}
\usepackage{xcolor}
\usepackage{booktabs}
\usepackage{url}

\def\BibTeX{{\rm B\kern-.05em{\sc i\kern-.025em b}\kern-.08em
    T\kern-.1667em\lower.7ex\hbox{E}\kern-.125emX}}

\begin{document}

\title{Where Poison Detectors Fail: A Comparative Failure and Latency\\
Analysis of RAG Defense Signals}

\author{\IEEEauthorblockN{1\textsuperscript{st} Tvisha Thakur}
\IEEEauthorblockA{\textit{School of Computer Science} \\
\textit{and Engineering} \\
\textit{Vellore Institute of Technology}\\
Vellore, India \\
tvisha.thakur2025@vitstudent.ac.in}
\and
\IEEEauthorblockN{2\textsuperscript{nd} Karishma}
\IEEEauthorblockA{\textit{School of Computer Science} \\
\textit{and Engineering} \\
\textit{Vellore Institute of Technology}\\
Vellore, India}
\and
\IEEEauthorblockN{3\textsuperscript{rd} Sree Vali}
\IEEEauthorblockA{\textit{School of Computer Science} \\
\textit{and Engineering} \\
\textit{Vellore Institute of Technology}\\
Vellore, India}
}

\maketitle

\begin{abstract}
Retrieval-Augmented Generation (RAG) pipelines trust whatever document a
retriever hands them, which makes knowledge poisoning, the insertion of a
malicious document into a trusted corpus, a practical and largely
undefended attack surface. We build a Poison Detection System (PDS) that
screens incoming documents using three published signals: instruction
pattern matching, embedding-based anomaly detection against a clean
baseline corpus, and language-model perplexity scoring. Rather than
proposing a new detector, our contribution is a comparative failure
analysis of how these signals combine and where they break. On a
77-document development set, naive equal-weight fusion reaches only
F1 = 0.79, not because any detector is weak, but because the embedding
detector's raw score range is far narrower than the other two, so simple
averaging dilutes it. Calibrating each detector to a common scale raises
F1 to 0.97, and a logistic-regression fusion on calibrated inputs reaches
0.98 and independently identifies the embedding detector as most
influential. A staged pipeline that defers the costliest detector matches
F1 = 0.97 with 70.5\% less detector latency. Red-teaming our own system
with ten evasive documents caught nine; the miss revealed a failure mode
we call \textit{boilerplate anchoring}, and a pattern added in response
brings the count to ten, though that patch is not independent evidence of
robustness. Two further results are negative: segment-level scoring did
not help, and on twelve benign documents from another domain the system
produced 12 false positives, all driven by the embedding detector,
showing that detector partly measures distance from its baseline domain
rather than malicious intent. All fusion numbers are development-set
results and are optimistic; a held-out evaluation protocol is under way.
\end{abstract}

\section{Introduction}

Retrieval-Augmented Generation systems answer a query by retrieving
documents from a corpus and conditioning a language model's response on
their contents. This architecture is now standard in production
assistants, but it inherits a quiet assumption: that every document in the
corpus is trustworthy. When that assumption fails, the result is
\textit{knowledge poisoning}, the deliberate insertion of a document
designed to manipulate downstream generation. This is not a hypothetical
concern; indirect prompt injection through externally sourced content has
been demonstrated against production LLM-integrated applications
\cite{yi2023benchmarking}, and dedicated RAG poisoning attacks that require
injecting only a single malicious passage per targeted query have since
been shown to be practical at scale \cite{zhang2025practical}.

A Poison Detection System (PDS) sits between the retriever and the
generator and screens each document before it is trusted. The individual
signals available for this task, keyword and instruction-pattern
matching, embedding-based anomaly detection, and language-model fluency
scoring, are each described in prior work. Our purpose is not to introduce
a fourth signal, but to ask: when these signals are combined, which ones
actually matter, why does a naive combination underperform, and where does
the resulting system still fail, including when it is tested on documents
unlike those it was built on.

\subsection{Contributions}
\begin{enumerate}
  \item A working three-detector PDS pipeline (pattern, embedding,
    perplexity) evaluated on a corpus of 55 real open-source README
    paragraphs and 60 constructed poisoned documents across four attack
    categories, plus 10 adversarial and 12 cross-domain documents.
  \item A diagnosed and corrected fusion failure: equal-weight averaging
    suppresses the most informative detector purely because of a
    raw-score-range mismatch, fixed by calibration and independently
    confirmed by a learned fusion model (F1 0.79 $\rightarrow$ 0.97 and
    0.98).
  \item A staged cascade reaching the same F1 as calibrated fusion with
    70.5\% less detector latency.
  \item An adversarial evaluation that names a specific weakness
    (boilerplate anchoring), a documented calibration-fragility analysis,
    and two honest negative results: segment-level scoring does not fix
    anchoring, and the embedding detector produced false positives on
    every benign document from a different domain.
\end{enumerate}

\section{Related Work}

\subsection{Indirect prompt injection and RAG poisoning}
Yi et al.\ introduce BIPIA, the first benchmark for indirect prompt
injection, and show that language models are broadly unable to
distinguish informational context from instructions embedded within it
\cite{yi2023benchmarking}. Hines et al.\ propose spotlighting, an
input-transformation defense that makes the provenance of untrusted text
explicit to the model, reducing injection success without requiring the
untrusted content to be filtered beforehand \cite{hines2024spotlighting}.
Su et al.\ evaluate RAG systems directly under adversarial poisoning and
find that a poisoned passage can influence both retrieval and generation,
motivating detection at the document level rather than relying on the
generator to notice an attack \cite{su2024robustrag}. Chen et al.\ study
whether indirect prompt injections can be both detected and removed once
retrieved, reinforcing that retrieved content should be treated as
untrusted input by default \cite{chen2025detectremove}. Zhang et al.\ show
that a poisoning attack requiring only a single injected document per
targeted query is practical and stealthy, the threat model we adopt
\cite{zhang2025practical}.

\subsection{Embedding-based anomaly detection}
Lee et al.\ propose CED, a training-free method for distinguishing
in-distribution from out-of-distribution text using embedding-space
differences, and note that normal and anomalous text can occupy
overlapping regions of embedding space, a difficulty directly relevant to
our own failure analysis \cite{lee2024ced}.

\subsection{Positioning of this work}
None of the three detection techniques used in PDS is original to this
paper. Our contribution is a comparative failure and latency analysis:
which signals earn their computational cost, how they should be combined,
and where the combination still breaks under adversarial pressure and under
domain shift.

% Threat model and formal problem statement (task 10.2)
% Grounded in: README.md, related_work.md, data/_provenance/, the attack
% categories already present in eval/results/evaluation_raw.csv (injection,
% authority, corpus, perplexity), and the adversarial/credential-exposure
% findings in eval/results/summary_report.md.

\section{Threat Model and Problem Statement}
\label{sec:threat-model}

\subsection{System Model}
We consider a Retrieval-Augmented Generation (RAG) pipeline consisting of a
document corpus $C$, a retriever $R$, and a downstream language model $M$.
At query time, $R$ selects a set of documents from $C$ and passes them to
$M$ as context. PDS sits between $R$ and $M$ (equivalently, it can be run
offline against $C$ itself): each candidate document $d$ is scored before
it is allowed to influence $M$'s output.

\subsection{Attacker Capabilities}
We assume an attacker who can introduce or modify at least one document
that is later retrieved into $C$ -- for example, by contributing content to
a corpus with open or loosely-moderated ingestion (a wiki, a shared
documentation repository, a user-submitted knowledge base), or by having
one of their pages indexed by a web-scraped source. Consistent with prior
work on practical RAG poisoning attacks requiring only a single injected
document rather than large-scale corpus control~\cite{zhang2025practical},
we do \emph{not} assume the attacker needs to control a large fraction of
the corpus: our adversarial evaluation set (\texttt{data/adversarial/})
specifically tests single-document, hand-crafted attacks.

The attacker's goal is to have $M$ follow an injected instruction, leak
information, or otherwise deviate from its intended behavior as a result of
processing the planted document.

\textbf{The attacker CANNOT:}
\begin{itemize}
  \item Modify the retriever's ranking algorithm or the language model's
    weights, prompts, or system instructions directly.
  \item Observe PDS's internal detector scores or decision boundaries at
    attack-construction time (black-box with respect to the deployed
    detector, though we do evaluate white-box-style adversarial attempts
    in \texttt{data/adversarial/}, where the attacker is assumed to know
    the \emph{general approach} -- pattern/embedding/perplexity scoring --
    but not the exact calibration bounds or thresholds).
  \item Guarantee retrieval: PDS assumes the defender controls the
    detection step, not whether a malicious document gets retrieved in the
    first place (that is the retriever's / corpus-governance's
    responsibility, out of scope here).
\end{itemize}

\subsection{Attack Surface / Categories}
Our evaluation dataset covers four primary categories, matching the labels
used throughout \texttt{eval/results/}:
\begin{enumerate}
  \item \textbf{Instruction injection} -- explicit override phrases
    (\emph{"ignore previous instructions"}, fake system tags).
  \item \textbf{Authority spoofing} -- content that impersonates a
    legitimate authority (a fake security policy, a fake administrator
    notice) to lend false credibility to an instruction.
  \item \textbf{Corpus poisoning} -- semantically out-of-place content
    inserted into an otherwise trusted corpus, without necessarily using
    explicit override language.
  \item \textbf{Perplexity artifacts} -- text that is unnaturally smooth
    (often a sign of adversarially optimized or template-generated attack
    text) or unnaturally garbled.
\end{enumerate}
Our adversarial red-teaming set adds a fifth, informally-named category
surfaced during evaluation: \textbf{credential exposure framed as
boilerplate} -- a malicious instruction (e.g., shipping default
credentials) phrased as neutral, generic documentation/license text
specifically to dilute embedding-anomaly and avoid pattern triggers. See
Section~\ref{sec:limitations} for how this was handled.

\subsection{Defender Assumptions}
PDS assumes the defender:
\begin{itemize}
  \item Has access to full document text before the trust decision is made
    (i.e., detection happens pre-retrieval-trust, not by inspecting $M$'s
    output after the fact).
  \item Has a small corpus of known-clean reference documents
    (\texttt{data/clean/}) to calibrate the embedding-anomaly detector
    against.
  \item Operates under a latency budget: detection must not make the RAG
    pipeline unusably slow, which motivates the staged/tiered cascade
    design (Section~\ref{sec:limitations} notes the latency measurement
    caveats).
  \item Does \emph{not} have ground-truth poisoned/clean labels at
    deployment time -- only at calibration time, using the fixed,
    pre-computed bounds described in \texttt{eval/calibrated\_fusion.py}.
\end{itemize}

\subsection{Formal Problem Statement}
For a document $d$, each detector $i \in \{\text{pattern}, \text{embedding},
\text{perplexity}\}$ produces a score $s_i(d) \in [0, 1]$ and a latency cost
$\ell_i(d)$. The fusion layer maps the score vector to a verdict
$v(d) \in \{\text{Allow}, \text{Flag for review}, \text{Reject}\}$.

We frame the staged pipeline as a sequential decision / early-exit cascade:
at each stage $i$, the defender either (a) stops and emits a verdict based
on scores collected so far, or (b) pays cost $\ell_{i+1}(d)$ to run the next
detector. The objective is to choose a stopping rule $\pi$ minimizing
expected total cost:
\[
  \mathbb{E}_{d \sim C}\Big[\, \alpha \cdot \mathbb{1}[\text{false negative}]
  + \beta \cdot \mathbb{1}[\text{false positive}]
  + \gamma \cdot \textstyle\sum_{i \in \pi(d)} \ell_i(d) \,\Big]
\]
where $\alpha, \beta, \gamma$ weight missed-detection cost, false-alarm
cost, and latency cost respectively, and $\pi(d)$ is the set of stages
actually run for document $d$ under policy $\pi$. The fixed-order staged
pipeline in \texttt{eval/staged\_pipeline.py} is a simple, hand-chosen
instance of $\pi$ (run cheap, high-precision detectors first; only pay for
perplexity when inconclusive); Phase 5's cost-aware cascade
(\texttt{fusion/cost\_aware.py}) is intended to choose $\pi$'s stage order
and early-exit thresholds more systematically.

\subsection{Explicitly Out of Scope}
\begin{itemize}
  \item Attacks on the retriever's ranking function itself (SEO-style
    manipulation to win retrieval, independent of document content).
  \item Attacks against the language model $M$ after a document has
    already been allowed through (jailbreaks that don't rely on retrieved
    content).
  \item Colluding multi-document attacks, where no single document is
    individually suspicious (planned future work, Phase 8.3).
  \item Adaptive attackers with white-box access to PDS's exact calibration
    bounds and thresholds (our adversarial set assumes attacker knowledge
    of the \emph{general} detection approach, not the specific deployed
    parameters).
\end{itemize}


\section{System Architecture}

\subsection{Detector A: instruction pattern matching}
A regular-expression detector scans document text for known override
phrases (variants of ``ignore previous instructions,'' fake
system-authority tags such as \texttt{[SYSTEM OVERRIDE]},
credential-exfiltration phrasing, and, after our red-teaming, phrasing that
exposes default or hard-coded credentials). Each pattern carries a fixed
weight; the document score is the maximum weight among matched patterns.
The detector is effectively free to run (about 0.1\,ms) but is blind to any
attack that does not use a known phrasing.

\subsection{Detector B: embedding anomaly detection}
A sentence-embedding model (all-MiniLM-L6-v2) encodes the document and a
reference corpus of trusted documents. The anomaly score is
$(1-\text{sim})/2$, where sim is the cosine similarity to the nearest
trusted document; a document far from every trusted example scores as more
anomalous. The detector costs about 75\,ms per document on our CPU.

\subsection{Detector C: perplexity/fluency scoring}
GPT-2 computes the perplexity of the document text. Very low perplexity
(unusually smooth, template-like) and very high perplexity (garbled) are
both treated as suspicious, with thresholds of 20 and 150. This is the most
expensive detector, at about 280\,ms per document on our CPU.

\subsection{Fusion strategies}
Let $s_i(d)\in[0,1]$ be detector $i$'s score on document $d$. The naive
approach is an unweighted average,
\begin{equation}
    s_{\text{fixed}}(d) = \frac{1}{3}\sum_{i=1}^{3} s_i(d).
\end{equation}
If detector $i$'s observed scores occupy $[\ell_i, h_i]$ and $h_i-\ell_i$
differs substantially across detectors, Eq.\ 1 under-represents the
detector with the narrowest range regardless of how discriminative it is.
We evaluate (1) naive fixed fusion; (2) calibrated fixed fusion, where each
detector's scores are min-max normalized to $[0,1]$ using bounds from the
evaluation data before averaging; (3) learned fusion, a logistic regression
trained via five-fold stratified cross-validation on raw and on calibrated
inputs (out-of-fold predictions reported); and (4) a staged pipeline in
which the cheapest detector runs first and each costlier detector runs only
if the earlier stages are inconclusive.

\section{Results and Discussion}
\label{sec:results}

\subsection{Dataset and protocol}
The corpus contains 55 clean documents, paragraphs from the READMEs of ten
open-source projects (Node.js, axios, Redis, FastAPI, pandas, Express,
Django, NumPy, requests, Flask), and 60 poisoned documents (15 per
category): direct instruction injection, fabricated system-authority
messages, on-topic ``corpus poisoning'' with harmful advice and no trigger
phrases, and garbled text. The embedding detector's reference corpus and
its evaluation set are disjoint: 70\% of clean documents (38) form the
trusted baseline, and the remaining 17, with all 60 poisoned documents,
are scored (77 total). Provenance for every document is logged. Thresholds
and calibration bounds were chosen on these same 77 documents, so all
numbers below are development-set results and are optimistic (see
Section~\ref{sec:limitations}).

\subsection{Fusion method comparison}
Table \ref{tab:fusion} summarizes the four fusion strategies.

\begin{table}[t]
\caption{Fusion method comparison (77-document development set)}
\label{tab:fusion}
\centering
\begin{tabular}{@{}lccc@{}}
\toprule
Method & Precision & Recall & F1 \\
\midrule
Naive fixed fusion & 1.00 & 0.65 & 0.79 \\
Learned fusion (raw scores) & 0.86 & 0.90 & 0.88 \\
Calibrated fixed fusion & 0.98 & 0.97 & 0.97 \\
Learned fusion (calibrated) & 0.97 & 0.98 & 0.98 \\
Staged pipeline & 0.94 & 1.00 & 0.97 \\
\bottomrule
\end{tabular}
\end{table}

The gap between naive and calibrated fusion is explained by score range.
The embedding detector's raw scores occupied $[0.180, 0.430]$, while the
pattern and perplexity detectors each spanned the full $[0,1]$ range.
Unweighted averaging therefore compresses the embedding detector by
construction, though a per-detector threshold sweep shows it is the
strongest signal (F1 = 0.92 alone at its best threshold, against 0.55 and
0.49 for pattern and perplexity). Calibrating each detector to $[0,1]$
before averaging raises F1 from 0.79 to 0.97 with no change to any
detector's scoring logic.

The learned fusion corroborates this. A logistic regression on raw scores
reaches only F1 = 0.88 and gives the embedding detector the smallest
coefficient of the three (1.09, against 1.82 and 1.51) despite its being
the most discriminative signal. Trained on calibrated inputs the same model
reaches F1 = 0.98 and ranks the embedding detector first (coefficient
2.67). A hand-derived calibration argument and a cross-validated model thus
converge on the same conclusion.

\subsection{Failure map by attack category}
Table \ref{tab:failuremap} reports, per category, the fraction of poisoned
documents caught by each detector at its own best threshold, and by
calibrated fusion at threshold 0.20.

\begin{table}[t]
\caption{Per-category detection rate}
\label{tab:failuremap}
\centering
\begin{tabular}{@{}lcccc@{}}
\toprule
Category & Pattern & Embed. & Perplex. & Fusion \\
\midrule
Injection & 67\% & 73\% & 33\% & 100\% \\
Authority & 73\% & 93\% & 0\% & 93\% \\
Corpus poisoning & 13\% & 100\% & 0\% & 93\% \\
Garbled/perplexity & 0\% & 100\% & 100\% & 100\% \\
\bottomrule
\end{tabular}
\end{table}

No single detector covers every category, the central argument for an
ensemble. Corpus poisoning, harmful advice in ordinary technical language
with no override phrases, is caught almost entirely by the embedding
detector. As the cross-domain result below shows, part of that success is
likely topic distance from the README baseline rather than detection of
malicious intent.

\subsection{A regex gap found and closed by failure analysis}
Inspection of documents missed by calibrated fusion identified two
injection documents containing ``ignore \textit{your} previous
instructions,'' which the original regular expression did not match because
it allowed only one optional modifier before ``instructions.'' Generalizing
the pattern to accept a possessive pronoun raised injection-category
detection from 87\% to 100\% and calibrated fusion F1 from 0.95 to 0.97,
verified by re-running the full corpus to confirm no new false positives.

\subsection{Staged pipeline and latency}
Running all three detectors on every document costs 357.4\,ms on average in
our latest run. The staged pipeline (pattern, then embedding, then
perplexity only if inconclusive) costs 105.3\,ms, a 70.5\% reduction, with
F1 = 0.97. The embedding stage ran on 70\% of documents and the perplexity
stage on 17\%. The most accurate single signal need not be the most
expensive to deploy, provided cheaper signals first resolve the documents
they can. This is a replay of latencies logged in one evaluation run; an
earlier run gave 72.5\%, so the figure is indicative rather than precise.

\subsection{Adversarial red-teaming}
We wrote ten documents meant to defeat all three detectors: each avoids
the original pattern trigger phrases, is fluent to avoid a high perplexity
score, and imitates README vocabulary to keep embedding distance small.
Calibration bounds were frozen from the original evaluation rather than
recomputed, to avoid circularity. In the first round, nine of ten were
caught. The evasion, \texttt{adv\_010}, opened with ``this project is
distributed under the MIT license'' and then described build artifacts
published with default demo credentials. We call this
\textit{boilerplate anchoring}: prefixing a malicious clause with text so
common in the reference domain that it pulls the whole-document embedding
toward the baseline and dilutes the payload (combined score 0.106 against
a threshold of 0.15). We did not find prior work describing this effect in
this setting.

We then added two credential-exposure patterns (default/preconfigured
credentials, hard-coded secrets). They match three of the ten adversarial
documents and none of the 55 clean or 12 cross-domain documents, and all
ten are now caught (\texttt{adv\_010}: 0.406). Because the patterns were
written after reading \texttt{adv\_010}'s wording, this is a fix for a
found gap, \textbf{not} independent evidence of robustness; the embedding
and perplexity detectors alone still miss that document.

\subsection{Calibration fragility}
\label{sec:calibration-fragility}
Min-max bounds for the embedding detector are set by one document at each
extreme. Removing the single document that sets a bound shifts the lower
bound from 0.180 to 0.205 and the whole dataset's normalized embedding
scores by 2.1\% on average, roughly 1.8x more than under 5th/95th-percentile
bounds (1.1\%). The pattern and perplexity detectors are unaffected, since
their scores tie at 0 and 1. More stable calibration is not automatically
better detection: under quantile bounds \texttt{adv\_010} receives a lower
embedding contribution (0.215) than under min-max (0.318), so a switch
would require re-tuning the threshold.

\subsection{Cross-domain false positives}
Twelve benign documents (HR, support and office policies) were scored with
the same baseline, frozen bounds and threshold. All 12 were flagged, each
driven mainly by the embedding detector. Their median embedding score
(0.37) is essentially that of the poisoned documents (0.36) and well above
that of in-domain clean documents (0.27); only 1 of 17 held-out in-domain
clean documents was flagged at threshold 0.20. The baseline contains only
software READMEs, so any document in a different style looks anomalous even
when harmless. The embedding detector therefore partly measures distance
from the baseline domain, not malicious intent, and the in-domain F1 of
0.97 must not be read as domain-general.

\subsection{Negative result: segment-level scoring}
To address boilerplate anchoring we scored each sentence independently with
three aggregations (maximum, a blend of maximum and mean, and a high-threshold
veto). None helped: maximum aggregation raised clean false positives from 4
to 9 of 17 at the original threshold, and after re-tuning the peak F1 values
were 0.90, 0.90 and 0.92, none above the 0.92 of the unmodified
document-level embedding detector. The highest segment score of
\texttt{adv\_010} (0.319) was below the maximum segment score of 5 of 17
clean documents. Sentence-level scores are noisier than the document-level
average, so the anchoring weakness is better addressed at the pattern level
or by a different calibration scheme than by finer-grained embeddings.

% Limitations and negative-results section (task 10.3)
% Every item below names which later phase/experiment is meant to address it,
% per the roadmap's own "Done when" criterion for this task.

\section{Limitations and Negative Results}
\label{sec:limitations}

\subsection{Dataset Scale and Composition}
The evaluation set reported in Section~\ref{sec:results} totals 77
documents (55 clean, 60 poisoned across 4 categories, partially
overlapping with the clean split) plus 10 adversarial and 12 cross-domain
documents, and a meaningful share of the clean and poisoned documents were
authored by the team rather than drawn from independent real-world
sources. Results at this scale should be read as a methodology
demonstration, not yet as evidence the detectors generalize broadly.
\emph{Addressed by:} Phase 2, which scales the dataset to 300+ documents
drawn from public sources (BIPIA, PoisonedRAG, deepset/prompt-injections,
Open-Prompt-Injection, HackAPrompt) with licenses and provenance logged
per source.

\subsection{No Comparison Against Published Baselines}
All F1 numbers in Section~\ref{sec:results} compare fusion variants of
PDS's own three detectors against each other -- there is currently no
comparison against existing published prompt-injection/poisoning
detectors (e.g., a DeBERTa-based classifier, Meta's Prompt Guard). We
therefore cannot yet claim PDS outperforms -- or even matches -- the
current state of the art; we can only claim that, among our own fusion
strategies, calibration provides a large, measurable improvement.
\emph{Addressed by:} Phase 3 (baseline selection, wrapping, and a
protocol-compliant comparison table).

\subsection{Negative Results: Alternative Fusion Strategies}
Beyond the fusion variants reported in the main results table, the team
also implemented and evaluated three additional fusion strategies --
segmented scoring, score blending, and a veto rule (any single detector
above a high threshold forces a Reject regardless of the others) -- none
of which improved on calibrated fixed fusion's F1 on the development set.
\textbf{[Note to team: exact F1/precision/recall figures for these three
variants are not yet on \texttt{main} as of this writing -- task 0.1 moves
these scripts from local to \texttt{eval/experiments/}. Fill in the actual
numbers here once that lands, rather than leaving this qualitative.]} We
report this as a negative result deliberately: it indicates that, on this
dataset, the dominant lever for accuracy was fixing the score-scale
mismatch (calibration) rather than further restructuring how the three
scores are combined.

\subsection{Calibration Fragility}
Section~\ref{sec:calibration-fragility} shows that min-max calibration
bounds for the embedding detector are set by a single document at each
extreme: removing just that one document shifts the whole dataset's
normalized scores by an average of 2.1\%, roughly 1.8x more than the
equivalent shift under quantile (5th/95th percentile) bounds. This means
the current, deployed calibration is more sensitive to which exact
documents happened to be in the one-time calibration set than it ideally
should be. Critically, the experiment also shows that quantile calibration
being more \emph{stable} does not automatically make it more
\emph{effective} at the current decision threshold -- under quantile
bounds, the one evasive adversarial document (\texttt{adv\_010}) actually
scores \emph{further} from the Reject threshold, not closer. Switching
calibration methods would require re-tuning the decision threshold
alongside it, not a drop-in swap.
\emph{Addressed by:} Phase 5's conformal calibration work, which chooses
thresholds with a formal false-positive-rate guarantee rather than an
ad-hoc fixed point, making this trade-off explicit and tunable instead of
incidental.

\subsection{Single Documented Evasion, Narrow Attacker Model}
Of 10 hand-crafted adversarial documents, 1 (\texttt{adv\_010.txt})
evaded detection, via credential-exposure content framed as generic
license/release boilerplate. All 10 adversarial documents, caught or not,
were constructed by the team with knowledge of the detection approach in
general terms; none used obfuscation techniques (homoglyphs, zero-width
characters, base64 encoding, sentence-splitting across documents) that a
more sophisticated attacker might employ. The current 90\% catch rate on
adversarial documents should not be read as an estimate of robustness
against an adaptive, obfuscation-aware attacker.
\emph{Addressed by:} Phase 8 (obfuscation and detector-aware attack
evaluation, paraphrase-robustness testing, multi-document and
position-sensitivity testing).

\subsection{Latency Measurement Methodology}
The reported 70.5\% latency reduction from the staged/tiered pipeline
(\texttt{eval/staged\_pipeline.py}) is computed by replaying the
per-detector latencies already logged during the single evaluation run in
\texttt{eval/results/evaluation\_raw.csv} -- it is not yet a systematic
benchmark (no warm-up isolation, no repeated runs, no median/p95
reporting across $\geq$100 runs per detector, no hardware printout). The
relative ordering (pattern $\ll$ embedding $\ll$ perplexity) is almost
certainly robust, but the exact 72.5\% figure should be treated as
indicative, not a rigorously measured benchmark result.
\emph{Addressed by:} Phase 1's latency benchmark
(\texttt{eval/latency\_benchmark.py}), designed specifically to fix this.

\subsection{Calibration Staleness Over Time}
The fixed min-max (or, going forward, quantile) calibration bounds are
computed once from the development set and reused for every subsequent
document, including in the adversarial evaluation
(\texttt{eval/adversarial\_test.py} explicitly freezes these bounds rather
than recomputing them, by design, to avoid circularity -- see that
script's docstring). Whether and how quickly these bounds go stale as the
distribution of incoming documents shifts over time in a real deployment
is untested.
\emph{Addressed by:} not yet scheduled in the current roadmap; flagged
here as an open question for future work beyond the current project scope.


% Ethics / responsible-release section (task 10.2)

\section{Ethics and Responsible Release}
\label{sec:ethics}

\subsection{Dual-Use Considerations}
A detection system for RAG poisoning is inherently dual-use: understanding
how detectors fail is also, necessarily, a partial recipe for evading them.
Our own adversarial red-teaming (Section~\ref{sec:results}) surfaced a real
evasion technique -- framing a malicious instruction as generic
documentation/license boilerplate to dilute the embedding-anomaly signal
(``boilerplate anchoring''). We report this finding because withholding it
would leave practitioners unaware of a real blind spot in a published
class of defenses, which is a worse outcome than disclosing it alongside
a discussion of mitigations (Section~\ref{sec:limitations}). This follows
the standard responsible-disclosure norm in security research: publish the
weakness together with what is known about addressing it, rather than
publishing neither or publishing the weakness alone.

\subsection{Data Provenance and Consent}
\begin{itemize}
  \item Clean baseline documents (\texttt{data/clean/}) are drawn from
    real, public open-source project READMEs; sources and licenses are
    logged in \texttt{data/\_provenance/}.
  \item Poisoned and adversarial documents (\texttt{data/poisoned/},
    \texttt{data/adversarial/}) were authored by the team specifically for
    this evaluation. They are not drawn from, nor do they reproduce, any
    real production attack, exploit, or incident.
  \item No real user data, private corpora, or personally identifying
    information is used anywhere in the dataset.
\end{itemize}

\subsection{Scope of Claims}
Consistent with \texttt{related\_work.md}, this project does not claim to
introduce a novel detection algorithm. Its contribution is a comparative
failure, blind-spot, and latency analysis of three published detection
signals (instruction-pattern matching, embedding anomaly, perplexity
scoring) and their fusion. We are explicit about this scope so the work is
not mistaken for -- or oversold as -- a new, patent-worthy detection
technique in isolation from calibration method (Phase 5's conformal
calibration work is the project's primary candidate for a genuinely novel,
disclosable technical contribution; see \texttt{docs/prior\_art.md}).

\subsection{Pre-Publication Review}
Per the project roadmap (task 11.1), this work will be reviewed by the
institution's IP/patent cell \emph{before} any public posting of code,
data, or the paper, to confirm what can safely be disclosed and when. No
component of this project has been publicly posted outside the private
team repository as of this writing.

\subsection{Intended Use}
PDS is intended to help practitioners evaluate and compare RAG-poisoning
defenses, not to serve as a how-to guide for evading them. Where we discuss
a successful evasion, we pair it with the specific gap it exposes
(Section~\ref{sec:limitations}) and the concrete next step planned to close
it (e.g., the credential-exposure pattern added in the dataset in response
to the \texttt{adv\_010} finding).


\section{Future Work}
The main gaps follow from the limitations. First, a held-out train/validation/test
protocol with bootstrap confidence intervals, so that thresholds are never
chosen on the data they are reported on. Second, a larger multi-domain
dataset from public sources (BIPIA, PoisonedRAG and others) and comparison
against published detectors. Third, a domain-diverse embedding baseline,
tested by holding out cross-domain documents, to separate topic distance
from malicious intent. Fourth, threshold selection with a formal
false-positive guarantee via conformal calibration, which also addresses
the calibration fragility above. Fifth, fresh adversarial documents with
obfuscation, paraphrase and multi-document attacks, and a stronger language
model for the perplexity signal.

\section{Conclusion}
We built and evaluated a three-signal Poison Detection System for RAG
pipelines and found that the most useful results were about combination
and failure, not any single detector's accuracy. A raw-score range mismatch
produced an 18-point F1 gap between naive (0.79) and calibrated (0.97)
fusion, confirmed independently by a learned fusion model, and a staged
pipeline kept that accuracy at 70.5\% lower latency. Red-teaming exposed a
named failure mode, boilerplate anchoring, which a targeted pattern closes
but which finer-grained embeddings did not. Most importantly for anyone
reusing these detectors, the embedding signal flagged every benign
out-of-domain document, so its in-domain accuracy overstates what it
measures. All numbers are from a small development set; the planned
held-out protocol is the necessary next step before any broader claim.

\begin{thebibliography}{00}

\bibitem{yi2023benchmarking} J. Yi, Y. Xie, B. Zhu, K. Hines, E. Kiciman, G. Sun, X. Xie, and F. Wu, ``Benchmarking and Defending Against Indirect Prompt Injection Attacks on Large Language Models,'' \textit{arXiv preprint arXiv:2312.14197}, 2023.

\bibitem{hines2024spotlighting} K. Hines, G. Lopez, M. Hall, F. Zarfati, Y. Zunger, and E. Kiciman, ``Defending Against Indirect Prompt Injection Attacks With Spotlighting,'' \textit{arXiv preprint arXiv:2403.14720}, 2024.

\bibitem{su2024robustrag} J. Su, S. Zhu, S. Chen, N. Zhang, and D. Wang, ``Towards More Robust Retrieval-Augmented Generation: Evaluating RAG Under Adversarial Poisoning Attacks,'' \textit{arXiv preprint arXiv:2412.16708}, 2024.

\bibitem{lee2024ced} H. Lee et al., ``CED: Comparing Embedding Differences for Detecting Out-of-Distribution and Hallucinated Text,'' in \textit{Findings of the Association for Computational Linguistics: EMNLP 2024}, 2024.

\bibitem{chen2025detectremove} Y. Chen et al., ``Can Indirect Prompt Injection Attacks Be Detected and Removed?,'' in \textit{Proceedings of ACL 2025}, 2025.

\bibitem{zhang2025practical} B. Zhang, Y. Chen, M. Fang, Z. Liu, L. Nie, T. Li, and Z. Liu, ``Practical Poisoning Attacks against Retrieval-Augmented Generation,'' \textit{arXiv preprint arXiv:2504.03957}, 2025.

\end{thebibliography}

\end{document}